\documentclass[11pt,a4paper]{article}
\usepackage[T1]{fontenc}
\usepackage[utf8]{inputenc}
\usepackage{amsmath,amssymb,xcolor}
\usepackage[margin=25mm]{geometry}
\usepackage[colorlinks=true,urlcolor=blue]{hyperref}
\definecolor{nexorablue}{HTML}{2459B5}
\setlength{\parindent}{0pt}
\setlength{\parskip}{7pt}
\setlength{\emergencystretch}{2em}
\renewcommand{\familydefault}{\sfdefault}
\begin{document}

{\small\bfseries\color{nexorablue} NEXORA CPGE}\par
{\small MPSI / NEX-LOG-02 / Corrigé}\par
{\LARGE\bfseries Témoin individuel ou coefficient commun}\par
\vspace{0.25em}\hrule\vspace{0.6em}

\section*{Première assertion}
Pour chaque $x\in E$, le réel $y=2/x$ convient. Le témoin dépend de $x$, ce que l’ordre des quantificateurs permet.

\section*{Seconde assertion}
Il faudrait simultanément $y=2$ pour $x=1$ et $y=1$ pour $x=2$. La négation est donc vraie : pour tout $y\in\mathbb R$, il existe $x\in E$ tel que $xy\ne2$.


\end{document}
